% Provenance (never printed): sources/book2/BRIEFS.md, g12-13. Trade from endowments, scale and variety; value chains;
% winners and losers inside a country; a tariff and a quota in a small open economy (proved: two triangles); the arguments for
% protection with their evidence; trade agreements, customs unions and the single market; debate on trade and wages. Real
% numbers: figdata/grade-12/13-trade-and-trade-policy/ (WDI weighted mean applied tariffs; Kestria invented). Ledger: Amiti et
% al. 2019, Fajgelbaum et al. 2020, Autor-Dorn-Hanson 2013, Fajgelbaum-Khandelwal 2016 (abstracts of the working papers), the
% WTO in brief, Wikipedia on the Leontief paradox and the single market (date only; no estimate of its effect was read). No
% container-port photograph. The 2019 US value of the WDI tariff series (13.8) is left out as inconsistent with its series.
\chapter{Trade beyond Ricardo, and Trade Policy}\label{ch:g12:trade-and-trade-policy}

In 1990 India's imports paid, on average, a tariff of 56\,\% of their value. In 2022 they paid 4.6\,\%. China's average fell from 32\,\% in 1992 to 2.2\,\%, Brazil's from 32\,\% in
1989 to 7.3\,\%. Within three decades, the large economies that had protected their producers most opened to imports, while the European Union's tariffs, already low,
fell to 1.3\,\%.
% figdata: tariff_key (ind1990 56.4, ind2022 4.6, chn1992 32.2, chn2022 2.2, bra1989 31.9, bra2022 7.3, euu2022 1.3)

This chapter widens the reasons for trade beyond the comparative advantage of \cref{ch:g10:ricardo-model}, then measures who gains and who loses, and what tariffs do.

\begin{openingchart}{Average tariff on imports, weighted by the value of imports, per cent, 1988--2022; the lines join the years with data.}{World Bank, World Development
Indicators (TM.TAX.MRCH.WM.AR.ZS)}
\begin{tikzpicture}
\begin{axis}[omtime, width=0.9\linewidth, height=6.0cm, xmin=1988, xmax=2022, ymin=0, ymax=60, ytick={0,10,...,60}, ylabel={per cent}, unbounded coords=discard]
\addplot[omDef, thick] table[x=year, y=IND] {figdata/grade-12/13-trade-and-trade-policy/out/tariffs.csv};
\addplot[omThm, thick] table[x=year, y=CHN] {figdata/grade-12/13-trade-and-trade-policy/out/tariffs.csv};
\addplot[omProp, thick] table[x=year, y=BRA] {figdata/grade-12/13-trade-and-trade-policy/out/tariffs.csv};
\addplot[omMeth, thick] table[x=year, y=MEX] {figdata/grade-12/13-trade-and-trade-policy/out/tariffs.csv};
\addplot[black, thick] table[x=year, y=KOR] {figdata/grade-12/13-trade-and-trade-policy/out/tariffs.csv};
\addplot[gray, thick] table[x=year, y=EUU] {figdata/grade-12/13-trade-and-trade-policy/out/tariffs.csv};
\node[font=\scriptsize, omDef, anchor=west] at (axis cs:1990.4,55) {India};
\node[font=\scriptsize, omThm, anchor=south west] at (axis cs:1992.5,31) {China};
\node[font=\scriptsize, omProp, anchor=west] at (axis cs:1989.5,25) {Brazil};
\node[font=\scriptsize, omMeth, anchor=south] at (axis cs:2003,16) {Mexico};
\node[font=\scriptsize, anchor=south] at (axis cs:2016,9) {South Korea};
\node[font=\scriptsize, gray, anchor=north] at (axis cs:2012,1.3) {European Union};
\end{axis}
\end{tikzpicture}
\end{openingchart}

\begin{recall}
From the grades 6--9 volume: exports, imports, tariffs, import quotas, and trade as a share of output. From \cref{ch:g10:ricardo-model}: comparative advantage and the \omterm{def:g10:ricardo-model:terms-of-trade}{terms of
trade}. From \cref{ch:g10:surplus-and-efficiency}: consumer and \omterm{def:g10:surplus-and-efficiency:producer-surplus}{producer surplus}. From \cref{ch:g10:tax-incidence}: who pays a tax. From \cref{ch:g11:public-choice-first-look}:
\omterm{def:g11:public-choice-first-look:rent-seeking}{rent seeking}.
\end{recall}

\section{Why countries trade}\label{sec:g12:trade-and-trade-policy:why}

\begin{definition}[Factor endowment]\label{def:g12:trade-and-trade-policy:endowment}
A country's \emph{factor endowment}\index{factor endowment} is the quantity it has of each factor of production: land, capital, skilled and unskilled labour.
\end{definition}

\begin{model}{The endowments model (Heckscher--Ohlin)}
\index{Heckscher--Ohlin model}%
\assumes Countries share the same technologies but differ in their endowments; goods differ in the factors they use intensively.
\predicts A country exports the goods that use its abundant factor intensively; trade raises the return to its abundant factor and lowers the return to its scarce factor.
\tested Mixed. The first test, in 1953, found American exports more labour-intensive than its imports, the reverse of the prediction for a capital-rich country; a
  later test across many countries found far less trade in factors than the model predicts, and fitted better once technology was allowed to differ across countries
  (Trefler 1995).
% ledger: b3fe:tre.missing
\breaks When technologies differ, which is Ricardo's reason for trade; and it says nothing of the large trade between similar countries.
\end{model}
% ledger: g12:leontief

\begin{definition}[Intra-industry trade]\label{def:g12:trade-and-trade-policy:intra-industry}
\emph{Intra-industry trade}\index{intra-industry trade} is the export and import of goods of the same kind, such as cars sold both ways between two countries.
\end{definition}

Similar countries trade cars for cars because each firm produces a few varieties at large scale, and buyers like variety. Trade lets each variety be made where its producer
is, for a larger market, at a lower cost per unit. Within any industry, the firms that export are the larger and more productive ones: American exporting firms in 2002 employed about two and a
half times as many workers as other firms of the same industry, and their value added per worker was about a tenth higher (Bernard, Jensen, Redding and Schott 2007);
% ledger: b3sv:bjrs.t3
 opening to trade moves workers from the
less productive firms, which shrink, to the more productive, which grow.

\begin{definition}[Global value chain]\label{def:g12:trade-and-trade-policy:value-chain}
A \emph{global value chain}\index{global value chain} is the sequence of stages, from design and parts to assembly and sale, by which a product is made in several countries.
\end{definition}

\begin{omfigure}
\begin{tikzpicture}[font=\scriptsize, node distance=0.4cm, box/.style={draw, rounded corners, fill=omDef!10, align=center, minimum height=1.0cm, text width=2.1cm}]
\node[box] (a) {design\\ \emph{Country A}};
\node[box, right=of a] (b) {chips and parts\\ \emph{Countries B, C}};
\node[box, right=of b] (c) {assembly\\ \emph{Country D}};
\node[box, right=of c, fill=omThm!10] (d) {sale\\ \emph{everywhere}};
\draw[->, thick] (a) -- (b); \draw[->, thick] (b) -- (c); \draw[->, thick] (c) -- (d);
\end{tikzpicture}

{\small A product made in stages across invented countries. Parts may cross borders several times, so a tariff on one stage raises the cost of every later one.}
\end{omfigure}

\section{Winners and losers inside a country}\label{sec:g12:trade-and-trade-policy:losers}

Trade raises a country's total income, as \cref{ch:g10:ricardo-model} proved, but not everyone's. Workers in industries that compete with imports can lose their jobs, and
those in the places where such industries gather can see wages fall for years.

\begin{debate}[Trade and wages]
% ledger: g12:adh2013, g12:fk2016, g12:leontief
\emph{Large local losses.} Autor, Dorn and Hanson (2013) compared American local labour markets more and less exposed to Chinese imports from 1990 to 2007: the more exposed had
higher unemployment, lower participation and lower wages, and more people on benefits; the rise in imports explains, conservatively, a quarter of the fall in American
manufacturing employment over those years. The design measures differences between places, not the national total, which also includes gains elsewhere.

\emph{Gains spread across buyers.} Fajgelbaum and Khandelwal (2016) estimate that trade typically favours poorer households, who spend more of their income on traded goods.
Their estimate rests on a model of demand fitted to aggregate data, not on observed households.

\emph{How large overall.} How much of the rise in the spread of wages trade explains, against technology, has no agreed estimate; the two studies above
measure different things, one the loss in exposed places, the other the gain to buyers, and neither gives the national share.
\end{debate}

\section{A tariff and a quota}\label{sec:g12:trade-and-trade-policy:tariff}

\begin{definition}[World price and small open economy]\label{def:g12:trade-and-trade-policy:world-price}
The \emph{world price}\index{world price} of a good is the price at which it is traded on world markets. A \emph{small open economy}\index{small open economy}%
\phantomsection\label{def:g12:trade-and-trade-policy:small-open-economy} trades at world prices that its own purchases and sales are too small to change.
\end{definition}

\begin{omfigure}
\begin{tikzpicture}
\begin{axis}[width=0.66\linewidth, height=6.0cm, xmin=0, xmax=90, ymin=20, ymax=80, xlabel={metres of cloth, millions}, ylabel={price, kests},
  xtick={20,28,52,60}, ytick={40,48}, grid=major, grid style={gray!20}, clip=false]
\fill[omThm!25] (axis cs:28,40) -- (axis cs:28,48) -- (axis cs:52,48) -- (axis cs:52,40) -- cycle;
\fill[omDef!35] (axis cs:20,40) -- (axis cs:28,48) -- (axis cs:28,40) -- cycle;
\fill[omDef!35] (axis cs:52,40) -- (axis cs:52,48) -- (axis cs:60,40) -- cycle;
\fill[omProp!25] (axis cs:0,40) -- (axis cs:20,40) -- (axis cs:28,48) -- (axis cs:0,48) -- cycle;
\addplot[omDef, thick, domain=20:80, samples=2] ({100 - x}, {x}) node[pos=0.95, right, font=\scriptsize] {demand};
\addplot[omThm, thick, domain=20:80, samples=2] ({x - 20}, {x}) node[pos=0.95, right, font=\scriptsize] {supply};
\addplot[black, thick, domain=0:90, samples=2] {40} node[pos=1, right, font=\scriptsize] {world price};
\addplot[black, dashed, domain=0:90, samples=2] {48} node[pos=1, right, font=\scriptsize] {with tariff};
\node[font=\scriptsize] at (axis cs:12,44) {a};
\node[font=\scriptsize] at (axis cs:26,42.5) {b};
\node[font=\scriptsize] at (axis cs:40,44) {c};
\node[font=\scriptsize] at (axis cs:54,42.5) {d};
\end{axis}
\end{tikzpicture}

{\small A 20\,\% tariff on cloth in Kestria, with the numbers of the weekend problem. The price rises from 40 to 48 kests; buyers lose $a + b + c + d$; producers gain $a$;
the government collects $c$ on the remaining imports; $b$ and $d$ are lost to all.}
\end{omfigure}

\begin{proposition}[The cost of a tariff]\label{prop:g12:trade-and-trade-policy:tariff}
In a \omterm{def:g12:trade-and-trade-policy:world-price}{small open economy} with linear demand and supply, a tariff $t$ per unit that leaves some imports causes a \omterm{def:g10:what-prices-do:dwl}{deadweight loss} of $\tfrac12 t\,\Delta Q_s + \tfrac12 t\,\Delta Q_d$,
where $\Delta Q_s$ is the rise in domestic production and $\Delta Q_d$ the fall in consumption.
\end{proposition}
\begin{proof}
The domestic price rises by $t$. \omterm{def:g10:surplus-and-efficiency:consumer-surplus}{Consumer surplus} falls by the area between the two prices left of the demand curve, a trapezium $a + b + c + d$. \omterm{def:g10:surplus-and-efficiency:producer-surplus}{Producer surplus} rises by
the trapezium $a$ left of the supply curve. Tariff revenue is $t$ times the new imports, the rectangle $c$. The net change is $-(b + d)$, two triangles of height $t$ and
bases $\Delta Q_s$ and $\Delta Q_d$.
\end{proof}

Triangle $b$ is the cost of producing at home, at a cost above the \omterm{def:g12:trade-and-trade-policy:world-price}{world price}, cloth that could have been imported; triangle $d$ is the value of the cloth that buyers no
longer buy, though they valued it above its \omterm{def:g12:trade-and-trade-policy:world-price}{world price}. An import quota that lets in the same quantity raises the price by the same amount, but the rectangle $c$ goes to
whoever holds the import licences, not to the government.

\begin{model}{A tariff in a small open economy}
\assumes A fixed \omterm{def:g12:trade-and-trade-policy:world-price}{world price} and competitive domestic markets.
\predicts The domestic price rises by the full tariff; buyers lose, domestic producers and the treasury gain, and the country loses two triangles.
\tested After the American tariffs of 2018 the prices of the targeted imports, before duty, did not fall: the tariffs passed fully into what American buyers paid.
\breaks When domestic sellers can set their prices: after the American tariff on washing machines in 2018, the price of washers rose by nearly 12\,\%, and the
  price of dryers, which bore no tariff, rose by as much (Flaaen, Hortaçsu and Tintelnot 2020). A large country's tariff can also lower the
  \omterm{def:g12:trade-and-trade-policy:world-price}{world price} of what it buys, shifting part of the cost to foreign sellers; the American record of 2018 found no such shift.
% ledger: g9:fht-2020-price-effect, g9:fht-2020-dryers
\end{model}
% ledger: g12:arw2019, g12:fgkk2020

\begin{record}[The American tariffs of 2018--19]
% ledger: g12:arw2019, g12:fgkk2020
The United States raised tariffs on many imports in 2018, and its trading partners retaliated. Both studies of the episode found complete pass-through to the
duty-inclusive prices paid by Americans and large falls in imports and exports. The losses to American consumers and firms buying imports were about 51 billion dollars, or
0.27\,\% of GDP; after counting the tariff revenue and the gains of domestic producers, the net loss of real income was 7.2 billion dollars, 0.04\,\% of GDP. By the end of
2018 the cost was running at about 1.4 billion dollars a month.
\end{record}

\section{The arguments for protection}\label{sec:g12:trade-and-trade-policy:arguments}

\begin{definition}[Infant-industry argument]\label{def:g12:trade-and-trade-policy:infant-industry}
The \emph{infant-industry argument}\index{infant-industry argument} holds that a young industry should be protected until it has learned to produce at the costs of its foreign
rivals.
\end{definition}

The argument holds only if the industry's future cost savings, in \omterm{def:g12:present-value:present-value}{present value}, exceed the cost of protecting it, and if lenders would not finance the learning
themselves. Protection, once granted, creates producers who gain from keeping it: an industry that never grows up can lobby to stay protected, the \omterm{def:g11:public-choice-first-look:rent-seeking}{rent seeking} of
\cref{ch:g11:public-choice-first-look}. A large country can gain from a tariff through better \omterm{def:g10:ricardo-model:terms-of-trade}{terms of trade}, but only if others do not retaliate. A country may also keep
the capacity to make goods it would need in a war, a cost paid for security.

\begin{definition}[Anti-dumping duty]\label{def:g12:trade-and-trade-policy:anti-dumping}
An \emph{anti-dumping duty}\index{anti-dumping duty} is a tariff on an import judged to be sold below its price at home or below its cost.
\end{definition}

\section{Trade agreements and the single market}\label{sec:g12:trade-and-trade-policy:agreements}

\begin{definition}[Trade agreement, customs union, single market]\label{def:g12:trade-and-trade-policy:trade-agreement}
A \emph{trade agreement}\index{trade agreement} is a treaty in which countries lower their barriers to each other's goods. A \emph{customs union}\index{customs
union}\phantomsection\label{def:g12:trade-and-trade-policy:customs-union} is an agreement that removes tariffs between its members and sets common tariffs on the rest of the
world. A \emph{single market}\index{single market}\phantomsection\label{def:g12:trade-and-trade-policy:single-market} goes further: goods, services, capital and people move
between its members under common rules, without checks at the borders.
\end{definition}

The World Trade Organization, which began on 1 January 1995, succeeded the General Agreement on Tariffs and Trade, under which tariffs had been negotiated down since 1947. The
European Union's \omterm{def:g12:trade-and-trade-policy:trade-agreement}{single market} took effect on 1 January 1993. An agreement that lowers tariffs only between its members can draw imports away from cheaper producers outside:
its gain from new trade must be set against that diversion.
% ledger: g12:wto1995, g12:esm1993

\section{Exercises}

\begin{exercise}[$\star$]\label{exo:g12:trade-and-trade-policy:1}
Demand is $Q_d = 100 - 2P$ and supply $Q_s = 2P - 20$; the \omterm{def:g12:trade-and-trade-policy:world-price}{world price} is 20. Find imports under free trade, and the domestic price and imports with a tariff of 5 per
unit.
\end{exercise}

\begin{exercise}[$\star$]\label{exo:g12:trade-and-trade-policy:2}
The \omterm{def:g12:trade-and-trade-policy:world-price}{world price} is 50 and a country imposes a 10\,\% tariff. What is the domestic price in a \omterm{def:g12:trade-and-trade-policy:world-price}{small open economy}?
\end{exercise}

\begin{exercise}[$\star$]\label{exo:g12:trade-and-trade-policy:3}
From the opening chart, give China's average tariff in 1992 and 2022.
\end{exercise}

\begin{exercise}[$\star$]\label{exo:g12:trade-and-trade-policy:4}
From the opening chart, which of the six had the highest average tariff in 2022? Compare it with India's in 1990.
\end{exercise}

\begin{exercise}[$\star$]\label{exo:g12:trade-and-trade-policy:5}
A quota of 20 units raises the domestic price from the \omterm{def:g12:trade-and-trade-policy:world-price}{world price} of 20 to 25. Find the licence rent.
\end{exercise}

\begin{exercise}[$\star\star$]\label{exo:g12:trade-and-trade-policy:6}
In \cref{exo:g12:trade-and-trade-policy:1}, find the tariff revenue and the \omterm{def:g10:what-prices-do:dwl}{deadweight loss}.
\end{exercise}

\begin{exercise}[$\star\star$]\label{exo:g12:trade-and-trade-policy:7}
In \cref{exo:g12:trade-and-trade-policy:1}, the tariff is 10. Find imports, revenue and the \omterm{def:g10:what-prices-do:dwl}{deadweight loss}.
\end{exercise}

\begin{exercise}[$\star\star$]\label{exo:g12:trade-and-trade-policy:8}
In \cref{exo:g12:trade-and-trade-policy:1}, compare the tariff of 5 with a quota of 20 units. What is the same, and who gets the rectangle?
\end{exercise}

\begin{exercise}[$\star\star$]\label{exo:g12:trade-and-trade-policy:9}
In \cref{exo:g12:trade-and-trade-policy:1}, the country is large: a tariff of 5 lowers the \omterm{def:g12:trade-and-trade-policy:world-price}{world price} from 20 to 18, so the domestic price is 23. Find imports, revenue, the
part of the revenue paid by foreign sellers, the \omterm{def:g10:what-prices-do:dwl}{deadweight loss} and the country's net gain or loss.
\end{exercise}

\begin{exercise}[$\star\star$]\label{exo:g12:trade-and-trade-policy:10}
A bicycle sells for 100 on the world market and uses 60 of imported parts, which pay no tariff. A 10\,\% tariff on bicycles raises their domestic price to 110. By how much
does it raise the value added that domestic assemblers can earn, in per cent? This is the effective rate of protection.
\end{exercise}

\begin{exercise}[$\star\star$]\label{exo:g12:trade-and-trade-policy:11}
In the 2018 record, the losses to buyers were 51 billion dollars and the net loss 7.2 billion. Explain the difference.
\end{exercise}

\begin{exercise}[$\star\star\star$]\label{exo:g12:trade-and-trade-policy:12}
For a large country, the model allows a tariff to lower the price foreign sellers receive. What did the record of 2018 show, and what does it imply about who paid?
\end{exercise}

\begin{exercise}[$\star\star\star$]\label{exo:g12:trade-and-trade-policy:13}
Protecting an industry costs 10 a year for five years. From year 6 on, it saves 8 a year for ever. At 5\,\%, compare the \omterm{def:g12:present-value:present-value}{present values}. What else must be true for protection
to be the right tool?
\end{exercise}

\begin{exercise}[$\star\star\star$]\label{exo:g12:trade-and-trade-policy:14}
A country rich in skilled labour opens to trade with a country rich in unskilled labour. According to the endowments model, whose wages fall in each country?
\end{exercise}

\begin{exercise}[$\star\star\star$]\label{exo:g12:trade-and-trade-policy:15}
An \omterm{def:g12:trade-and-trade-policy:anti-dumping}{anti-dumping duty} can be requested by domestic producers. Why might it be used more than the cases of true dumping justify?
\end{exercise}

\section{Problem: Kestria's Tariff on Cloth}

\begin{problem}[{Weekend problem --- a tariff, a quota and retaliation}]
\label{pb:g12:trade-and-trade-policy:1}
In Kestria, demand for cloth is $Q_d = 100 - P$ and domestic supply is $Q_s = P - 20$, in millions of metres, with $P$ in kests per metre. Kestria is small and can import any
quantity at the \omterm{def:g12:trade-and-trade-policy:world-price}{world price} of 40 kests.

\medskip
\noindent\textbf{Part I --- Free trade and a tariff.}
\begin{enumerate}
\item Find consumption, production and imports under free trade.
\item Kestria imposes a tariff of 20\,\%. Find the domestic price.
\item Find consumption, production and imports with the tariff.
\item Find the loss of \omterm{def:g10:surplus-and-efficiency:consumer-surplus}{consumer surplus} and the gain of \omterm{def:g10:surplus-and-efficiency:producer-surplus}{producer surplus}, in millions of kests.
\item Find the tariff revenue.
\end{enumerate}

\medskip
\noindent\textbf{Part II --- The deadweight loss.}
\begin{enumerate}[resume]
\item Find the two triangles and the \omterm{def:g10:what-prices-do:dwl}{deadweight loss}.
\item Check that the loss of consumers equals the gains of producers and of the government plus the \omterm{def:g10:what-prices-do:dwl}{deadweight loss}.
\item Kestria uses a quota of 24 million metres instead. What is the domestic price?
\item Find the licence rent. Who gets it if licences are sold at auction? If they are given to foreign exporters?
\item What is Kestria's total loss in the second case?
\end{enumerate}

\medskip
\noindent\textbf{Part III --- Retaliation.}
\begin{enumerate}[resume]
\item Brenland answers with a tariff on Kestria's grain, which lowers the price Kestria's grain exporters receive from 30 to 27 kests. Kestria's grain supply is $Q_s = 2P -
  20$ and demand $Q_d = 80 - 2P$. Find grain exports before and after.
\item Find the loss of grain producers.
\item Find the gain of grain buyers in Kestria.
\item Find Kestria's net loss on grain.
\item Find Kestria's total net loss, cloth and grain together.
\end{enumerate}

\medskip
\noindent\textbf{Part IV --- Interpretation.}
\begin{enumerate}[resume]
\item Who in Kestria gains from the cloth tariff? Who loses from the retaliation?
\item Why might cloth producers lobby for the tariff although Kestria as a whole loses?
\item Which result of the 2018 record does Part I resemble?
\item What would change if Kestria were large enough to lower the \omterm{def:g12:trade-and-trade-policy:world-price}{world price} of cloth?
\item State the result of the problem: the \omterm{def:g10:what-prices-do:dwl}{deadweight loss} of the tariff, and the net loss to Kestria after retaliation.
\end{enumerate}
\end{problem}
